\documentclass[10pt,a4paper]{amsart}
\usepackage[margin=19mm]{geometry}
\usepackage{amsmath,amssymb,mathtools}
\usepackage[scr=rsfs,cal=euler]{mathalfa}
\usepackage{xfrac}
\usepackage[unicode,hidelinks,bookmarks=false]{hyperref}
\newcommand{\ie}{i.e.\ }
\newcommand{\cf}{cf.\ }
\newcommand{\resp}{respectively\ }
\newcommand{\Subsection}{\subsection}
\providecommand{\nobreakdash}{\nobreak\hbox{-}}
\setlength{\emergencystretch}{2em}
\multlinegap0pt
\usepackage{bm}
\usepackage{bbm}
\usepackage{dsfont}
\usepackage{xspace}
\usepackage{cancel}
\usepackage{mathtools}
\usepackage{amsthm}
\makeatletter
\@ifundefined{lemma}{\newtheorem{lemma}{Lemma}}{}
\@ifundefined{theorem}{\newtheorem{theorem}{Theorem}}{}
\makeatother
\makeatletter
\def\FF{\mathbb{F}}
\def\NN{\mathbb{N}}
\expandafter\ifx\csname assumptionprime\endcsname\relax
\newenvironment{assumptionprime}[1]
  {\renewcommand{\theassumption}{\ref{#1}$'$}%
   \addtocounter{assumption}{-1}%
   \begin{assumption}}
\fi
\def\eqref#1{equation~(\ref{#1})}
\makeatother
\title[On A Necessary Condition For Posterior Inconsistency: New In]{On A Necessary Condition For Posterior Inconsistency: New Insights From A Classic Counterexample}
\author{Nicola Bariletto, Stephen G. Walker}
\date{Source: arXiv:2510.18126v1 (CC BY 4.0)}
\begin{document}
\maketitle

\noindent\emph{Source and licence.} On A Necessary Condition For Posterior Inconsistency: New Insights From A Classic Counterexample. Version arXiv:2510.18126v1 (CC BY 4.0), distributed under the Creative Commons Attribution 4.0 International licence, as recorded in the arXiv licence metadata for this exact version and shown on the abstract page https://arxiv.org/abs/2510.18126v1. Full rights notice: licenses/arxiv-2510.18126-CC-BY-4.0.txt.\par\smallskip

\noindent\textit{Editorial scope.} Counted unit: the Theorem whose source label is \texttt{thm:alpha\_beta\_inconsistency} in the article (source0.tex lines 993--1008, source label \texttt{thm:alpha\_beta\_inconsistency}), delivered together with the article's own complete original proof of it (source0.tex lines 1010--1066, one \texttt{proof} environment of 57 lines). No mathematics is written, repaired or re-derived: statement and proof are the article's own text line for line. Required context retained uncounted: lem:schwartz (lines 778--780). Every definition, display and label of those lines is retained unchanged. Omitted from this package and not redistributed: 1--777, 781--992, 1009--1009, 1067--1144. Nothing inside a retained excerpt is omitted or altered: every retained body is delivered exactly as the article's lines are written, so no omission receipt of any kind accompanies this package. Citations: the retained excerpts carry no citation command at all, so no bibliography entry of the article is transcribed and no reference is redistributed. Reference adaptation: none was needed and none was made; every reference inside the delivered unit resolves inside it, so no unresolved reference or citation is produced. Printed slips kept literal: no printed word, symbol, hypothesis or number of the counted statement or of its proof was altered. Layout and class: the article's own class is \texttt{article} with packages including amsmath, amsfonts, bm, inputenc, fontenc, booktabs, nicefrac, microtype, subfigure, epsf, epsfig, fancyhdr, graphics, graphicx; the wrapper is the lane's pinned \texttt{amsart} editorial wrapper at 10pt on A4, which restates the theorem-like environments the retained text needs and adds the article's own macro definitions that the retained text needs (\texttt{\textbackslash FF}, \texttt{\textbackslash NN}, \texttt{\textbackslash eqref}), with extra wrapper packages (bm, bbm, dsfont, xspace, cancel, mathtools). The delivered numbering is the unit's own. Assets: no retained span contains a figure environment, a graphics-inclusion command, a PDF-page inclusion, a table or a TikZ picture; no image file is copied into this package, the package declares no asset dependency and no image file is redistributed with it. Licence: version arXiv:2510.18126v1 of this article is distributed under the Creative Commons Attribution 4.0 International licence, as recorded in the arXiv licence metadata for this exact version and shown on the abstract page https://arxiv.org/abs/2510.18126v1. A full rights notice is delivered at licenses/arxiv-2510.18126-CC-BY-4.0.txt. Characters: every retained excerpt is pure ASCII, so the delivered engine, XeLaTeX with its default Latin Modern text font, reports no missing character.
\begin{lemma}\label{lem:schwartz}
    Assume that $f_\star$ lies in the KL support of $\Pi$. Then, for any weak neighborhood $U$ of $F_\star$, $\lim_{n\to\infty} \Pi(U^c \mid X_{1:n}) = 0$ almost surely. Moreover, $\int_\FF R_n(f)\,\Pi(\mathrm df) \geq e^{-\tau n}$ ultimately almost surely, for all $\tau > 0$.
\end{lemma}

\begin{theorem}\label{thm:alpha_beta_inconsistency}
    % Assume $f_\star$ is in the KL support of $\Pi$ and that, for some $0<\alpha < \beta<\infty$, the posterior is $(\alpha,\beta)$-inconsistent. Then, for any sufficiently small $\chi>0$, there exist $0<\xi<\chi$ and $\gamma_\chi\in[\alpha,\beta]$ such that
    % %
    % \begin{equation*}
    %     \gamma_\chi - \xi -\tau_1 \leq - n^{-1} \ln\Pi\left(\left\{e^{(\gamma_\chi -\xi)n}\leq R_n(f) \leq e^{(\gamma_\chi + \xi)n}\right\}\right) \leq \gamma_\chi + \xi + \tau_2
    % \end{equation*}
    % %
    % infinitely often with positive probability under $F_\star^\infty$, for any $\tau_1, \tau_2>0$.
    Assume that $f_\star$ is in the KL support of $\Pi$ and that, for some $0<\alpha < \beta<\infty$, the posterior is $(\alpha,\beta)$-inconsistent. Then there exists $\gamma\in[\alpha,\beta]$ such that, for all small $\chi>0$, one can find $\mu_1,\mu_2\geq 0$ with $\mu_1+\mu_2\leq \chi$ and
    %
    \begin{equation}\label{eq:theorem_alpha_beta_inconsistency}
        \gamma - \mu_1 - \tau_1 \leq -n^{-1}\ln\Pi\big(\big\{f\in\FF : \gamma-\mu_1 \leq n^{-1}\ln R_n(f)\leq \gamma + \mu_2\big\}\big) \leq \gamma + \mu_2 + \tau_2
    \end{equation}
    %
    infinitely often with positive probability, for all $\tau_1,\tau_2>0$.
\end{theorem}

\begin{proof}
    % For some $\delta>0$, we have
    % %
    % \begin{align*}
    %     \delta & < \Pi(\{f\in\FF : e^{\alpha n}\leq R_n(f)\leq e^{\beta n}\}\mid X_{1:n}) \\
    %     & \leq   \Pi(\{f\in\FF : e^{\alpha n}\leq R_n(f)\leq e^{cn}\}\mid X_{1:n}) \, + \, \Pi(\{f\in\FF : e^{cn}\leq R_n(f)\leq e^{\beta n}\}\mid X_{1:n})
    % \end{align*}
    % %
    % infinitely often with positive probability, where $c:=(\alpha+\beta)/2$. Therefore, at least one addendum must be greater than some positive constant infinitely often with positive probability. One can iterate this argument $\lceil\log_2((\alpha-\beta)/2\chi)\rceil$ times to find an interval of length $2\xi\leq 2\chi$, say $[\gamma_\chi - \xi, \gamma_\chi+\xi]$ for some $\gamma_\chi\in[\alpha,\beta]$, such that
    % %
    % \begin{equation*}
    %     \Pi(\{f\in\FF : e^{(\gamma_\chi - \xi) n}\leq R_n(f)\leq e^{(\gamma_\chi+\xi) n}\}\mid X_{1:n}) > \delta_\xi
    % \end{equation*}
    % %
    % infinitely often with positive probability, for some $\delta_\xi>0$.
    Without loss of generality, pick $\chi=1/k$ for $k\in\NN$ large enough. For some $\delta>0$, we have
    %
    \begin{align*}
        \delta & < \Pi(\{f\in\FF : e^{\alpha n}\leq R_n(f)\leq e^{\beta n}\}\mid X_{1:n}) \\
        & \leq   \Pi(\{f\in\FF : e^{\alpha n}\leq R_n(f)\leq e^{cn}\}\mid X_{1:n}) \, + \, \Pi(\{f\in\FF : e^{cn}\leq R_n(f)\leq e^{\beta n}\}\mid X_{1:n})
    \end{align*}
    %
    infinitely often with positive probability, where $c:=(\alpha+\beta)/2$. Therefore, at least one addendum must be greater than some positive constant infinitely often with positive probability. One can iterate this argument $\lceil\log_2((\alpha-\beta)/\chi)\rceil$ times to find an interval of length $\xi_\chi\leq \chi$, say $[\gamma_\chi - \xi_\chi/2,\, \gamma_\chi+\xi_\chi/2]$ for some $\gamma_\chi\in[\alpha,\beta]$, such that
    %
    \begin{equation*}
        \Pi(\{f\in\FF : e^{(\gamma_\chi - \xi_\chi/2) n}\leq R_n(f)\leq e^{(\gamma_\chi+\xi_\chi/2) n}\}\mid X_{1:n}) > \delta_\chi
    \end{equation*}
    %
    infinitely often with positive probability, for some $\delta_\chi>0$.

    When $B_{\chi,n}:=\big\{f\in\FF : e^{(\gamma_\chi - \xi_\chi/2) n}\leq R_n(f)\leq e^{(\gamma_\chi+\xi_\chi/2) n}\big\}$, because Lemma~\ref{lem:schwartz} yields $\int_\mathbb F R_n(f)\, \Pi(\mathrm df)\geq e^{-\tau n}$ ultimately almost surely for all $\tau>0$, we obtain
    %
    \begin{equation*}
        \int_{B_{\chi,n}} R_n(f)\, \Pi(\mathrm d f)\geq e^{-\tau n}
    \end{equation*}
    %
    infinitely often with positive probability for all $\tau>0$. Moreover, by a simple application of Markov's inequality and the first Borel-Cantelli lemma, we have
    %
    \begin{equation*}
        \int_{B_{\chi,n}} R_n(f)\, \Pi(\mathrm d f)\leq \int_{\mathbb F} R_n(f)\, \Pi(\mathrm d f)\leq e^{\tau n}
    \end{equation*}
    %
    ultimately almost surely for all $\tau>0$. By definition of $B_{\chi,n}$, we also have that
    %
    \begin{equation*}
        e^{(\gamma_\chi - \xi_\chi/2)n}\, \Pi(B_{\chi,n})\leq \int_{B_{\chi,n}} R_n(f)\, \Pi(\mathrm d f)\leq e^{(\gamma_\chi + \xi_\chi/2)n}\, \Pi(B_{\chi,n}).
    \end{equation*}
    %
    Therefore, combining all of the above yields
    %
    \begin{equation*}
        \gamma_\chi - \xi_\chi/2 - \tau_1 \leq -n^{-1}\ln\Pi(B_{\chi,n}) \leq \gamma_\chi + \xi_\chi/2 + \tau_2
    \end{equation*}
    %
    infinitely often with positive probability, for all $\tau_1, \tau_2>0$. Now  $(\gamma_{1/k})_{k\in\NN}$ is a Cauchy sequence contained in the closed set $[\alpha,\beta]$, so it has a limit $\gamma\in[\alpha,\beta]$.
    Defining $\mu_1:=\gamma - (\gamma_\chi-\xi_\chi/2)$ and $\mu_2:=(\gamma_\chi+\xi_\chi/2)-\gamma$, condition \eqref{eq:theorem_alpha_beta_inconsistency} follows.
\end{proof}

\end{document}
